\section{Introduction}
\label{sec:intro}

Late-interaction retrieval models such as ColBERT~\cite{khattab2020colbert,santhanam2022colbertv2} keep one vector per token.
They score a query $q$ against a document $d$ with MaxSim: each query token takes its best similarity over the document tokens, and the score is the mean of these maxima.
Many systems use MaxSim as a reranking stage: a first stage (BM25 or a single-vector index) returns some tens or hundreds of candidates, and MaxSim orders them.

Rerank calls are small.
In the coding-agent workload of this paper, the median call has \nm{wl.agent-ho.q50} query tokens and \nm{wl.agent-ho.pool50} candidates of about \nm{wl.agent-ho.dt50} tokens each (Section~\ref{sec:workload}).
At this size, a call takes tens of microseconds.
Fixed costs such as library dispatch, thread wake-up and input conversion can then matter as much as arithmetic.
Common tools, however, are built for large matrix products.
So which CPU implementation is fastest for these calls, at what quality, and under which conditions, is an empirical question.
Earlier work timed CPU MaxSim libraries as baselines for faster approximate methods~\cite{pony2026colbandit}, or for a single library~\cite{clavie2025maxsimcpu}. We found no public benchmark that compares CPU MaxSim implementations on captured rerank calls, with each implementation checked against its own scoring rule.

This paper contributes four things.
\begin{enumerate}
\item \textbf{A benchmark} (\benchname): a replay format for rerank calls; four workload families (captured agent calls, GenCodeSearchNet calls, a size sweep, and BEIR SciFact and NFCorpus with two public ColBERT models); adapters for eight implementations; and per-call timing records for every repetition.
\item \textbf{Scoring contracts.} The tested libraries compute different scores: they differ in normalization, in where an epsilon goes, in a $\max(0,\cdot)$ clamp, and in padding. We define each rule once, compute a float64 reference per rule, and check every arm against the reference of its own rule (Section~\ref{sec:contracts}). Independent hand-built fixtures test each rule.
\item \textbf{A measurement protocol} that keeps development and held-out data apart. All configurations, the repetition count, the quality margin and the analysis code were frozen before any held-out run (Section~\ref{sec:method}). We report two uncertainty numbers apart: a workload interval and an execution spread.
\item \textbf{Results on five CPUs} (Granite Rapids, Zen~5, Graviton4, Apple M3 Max, and Zen~3 \TODO{epyc}): speed under warm, idle-gap, one-thread and parallel-caller conditions; memory per stored token; and quality per latency budget. Every case where a rival is faster is in the main text (Section~\ref{sec:losses}).
\end{enumerate}

One implementation is an int4 kernel from the open-source code-search engine sweet-search (\sweet below).
It is careful engineering, not a new algorithm: SIMD register blocking, int4 decoding inside the kernel, ragged batches without padding, and its own worker pool (Section~\ref{sec:kernel}).
The authors of this paper also wrote that kernel.
We therefore designed the protocol so that its results do not depend on our judgment after the data were seen: the rival configurations come from a frozen shortlist chosen on development data, and the claims follow fixed rules.

\paragraph{What we claim.}
On the tested CPUs, this implementation reduced warm scoring latency for captured reranking workloads, relative to the tested library configurations, with exhaustive scoring over stored int4 representations and measured floating-point error.
We do not claim the fastest MaxSim in general, any end-to-end search speed-up, or results for GPUs.
A stronger claim, that fixed costs dominate at real call sizes, needs a cost decomposition that we did not run (Section~\ref{sec:limits}).

\paragraph{What we found besides speed.}
The reference checks found two new defects in other libraries and confirmed a known one (Section~\ref{sec:correctness}).
The quality test found that the production int4 scoring rule loses quality with one of the two public models, and a post-hoc analysis locates the cause in the normalization step (Section~\ref{sec:quality}).
